\chapter*{Sumário executivo}

Esta \textbf{Proposta} foi elaborada em cumprimento dos objetivos 1.1, 2 e 8 definidos para a Comissão de Comunicação Institucional

\section*{Estrutura geral da \emph{Política}}

A presente \textbf{Proposta de \emph{Política de comunicação institucional} para o Conselho Regional de Psicologia de São Paulo} está dividida em quatro capítulos e dois apêndices.

No primeiro, introdutório, são apresentados os conceitos e princípios que devem nortear a comunicação institucional do CRP~SP.

No segundo, ``Estrutura e organização'', o organograma da área de comunicação do Conselho é descrito com base nas normativas correspondentes.

No terceiro, ``Canais, serviços e produtos'', são definidos os canais oficiais e o portfólio de produtos de comunicação da autarquia.

O quarto, ``Serviços de comunicação e como solicitá-los'', apresenta os serviços prestados pela Coordenação de Comunicação e define quem pode solicitá-los e como essa solicitação deve ser feita.

O apêndice A, ``Redação oficial: normas gerais'', propõe a padronização da grafia do nome do Conselho Regional de Psicologia de São Paulo e de suas unidades.

Por fim, no apêndice B, ``Diretrizes para a comunicação de conselheiras e trabalhadoras do CRP~SP nas redes sociais'', são apresentadas orientações gerais e dicas para que a comunicação \emph{on-line} de todas as pessoas ligadas ao Conselho esteja adequada às normas que regem a administração pública.